
As agentic AI systems move beyond static question answering into open-ended, tool-augmented, and multi-step real-world workflows---including retrieval-grounded assistants that act on the documents and structured data they fetch---their increased authority poses greater risks of system misuse and operational failures. However, current evaluation practices remain fragmented, measuring isolated capabilities such as coding, hallucination, jailbreak resistance, or tool use in narrowly defined settings, with the term ``trustworthiness'' frequently invoked but rarely defined operationally. We argue that the central limitation is twofold: (i)~the absence of a concrete, measurable specification of \emph{what} agent trustworthiness means, and (ii)~the lack of a principled notion of \emph{representativeness} that would allow an agent's behaviour to be assessed over a socio-technical scenario distribution rather than a collection of disconnected benchmark instances. We address (i) by defining agentic trustworthiness as a five-property profile---Reliability, Robustness, Safety, Social-Ethical Alignment, and Operational Integrity---grounded in current AI risk frameworks; and we address (ii) by proposing the \textbf{Holographic Agent Assessment Framework (HAAF)}, a systematic paradigm that characterises this profile over a \emph{scenario manifold} spanning task types, tool interfaces, interaction dynamics, social contexts, and risk levels. HAAF integrates four complementary components: (i)~\emph{static cognitive and policy analysis}, (ii)~\emph{interactive sandbox simulation}, (iii)~\emph{social-ethical alignment assessment}, and (iv)~a \emph{distribution-aware representative sampling engine} that jointly optimises coverage and risk sensitivity---particularly for rare but high-consequence tail risks that conventional benchmarks systematically overlook. These components are connected through an iterative \emph{Trustworthy Optimization Factory} that converts red-team diagnoses into targeted blue-team interventions until deployment-readiness thresholds are met. We instantiate HAAF on a retrieval-grounded agentic sandbox in three steps: (a)~a single-system Factory cycle that reduces violation rate from 16.7\% to 4.2\%; (b)~a \emph{cross-model trustworthiness profile} over 13 contemporary agentic systems spanning seven model families (Llama, Mistral, Kimi, GLM, Qwen, GPT, DeepSeek) that surfaces property-level trade-offs no scalar leaderboard can capture; and (c)~a \emph{cross-family Factory generalisation experiment} in which the focal-model interventions are applied uniformly---without per-model or per-scenario tuning---to all 13 systems on the 100-scenario suite: all 13 systems improve, with five reducing risk by $\Delta\mathrm{RWF}\!>\!0.27$ and two reaching a perfect $\mathrm{RWF}\!=\!0.000$ (GPT-oss-120B: $0.481\!\to\!0.000$; Llama-3.1-8B: $0.216\!\to\!0.000$), while one intervention-resistant case (Qwen3-32B, $\Delta\!=\!+0.019$) is identified by the profile as a candidate for a fundamentally different second Factory cycle---direct evidence that property-targeted hardening generalises across model families while preserving visibility into where it does not. Code and data are available at \url{https://github.com/TonyQJH/haaf-pilot}.
